\chapter{Introduction}
\label{chap:intro}


In 2016, ProPublica \cite{angwin2016machine} published their report on racial biases in the COMPAS Recidivism Algorithm used by courts in the United States, finding that black defendants were predicted almost twice as likely to re-offend as their white counterparts when controlling for all other variables. This report among many others motivated the field of AI ethics to study demographic biases in algorithms and found such biases at scale \cite{Buolamwini2018GenderSI}. Importantly, these biases not only apply to simpler machine learning algorithms such as used by COMPAS, but also persist into today's Large Language Models (LLMs), even after explicit training to make models harmless \cite{hofmann2024ai-e4f}.

However, demographics-driven harm is not only about misrepresenting certain demographic groups in a harmful way, but also about systematically failing to protect them. A premier case for this is Yong et al.'s study finding that LLMs might relatively robustly refuse harm-inducing user requests if the prompt is in English, but fail to do so in 79\% of cases when the prompt is given in a low-resource language like Zulu \cite{yong2023low-resource-ec3}. Similarly, in the case of false-refusal, prior work has found that given certain demographics-loaded names, models would overrefuse less often for white men than other demographic groups \cite{li-etal-2024-chatgpt-doesnt, plaza-del-arco-etal-2025-yes}.

Evidence for this failure of protection based on demographics is up to this point only anecdotal and understudied in comparison to representational harm. A reason for this might be that AI safety, which deals primarily with safeguards in AI, is often seen in separation from AI ethics and not in a continuum.

In fact, as Lazar and Nelson \cite{lazar2023ai-a8d} criticised, defining \textit{AI Safety} as \textit{preventing and mitigating large scale risk from AI} is in itself value-laden and far from neutral. Accumulation of small-scale harm for many users may equally result in high societal risks legitimising everyday-safety-for-all as core safety, not peripheral \cite{gyevnar2025ai-cae}. From the methodological point of view, safety evaluations and methods may thus suffer from construct-invalidity and a coverage gap \cite{raji2021everything, rottger2025safetyprompts}.

To the best of our knowledge, however, this failure of protection based on demographic coverage has not been studied systematically, raising the question to what extent this is a systematic problem, which we coin "safety discrimination" in this work:\\

\noindent \textit{\textbf{Definition Safety Discrimination:}  When a person, based on their demographic background, is deprived of the privilege of safeguards that protect them from harmful behaviour and output of an AI.}\\

Safety discrimination could further be seen from two perspectives: first – \textit{methodological safety discrimination} – are the methods we develop to evaluate and mitigate risks designed in a way that includes users from diverse demographics, or are we just relying on the hope that methods generalise to all users? Generalisability may be a false premise \cite{raji2021everything} under the lens of evaluations as measurements of a well-defined construct \cite{wallach2025position}. But beyond a well-defined construct, safety is itself context-dependent, and the same AI response can be safe or very unsafe depending on who the user is \cite{kempermann2026safewhomrethinkingevaluate, Weidinger2023SociotechnicalSE}.

Second –  \textit{behavioural safety discrimination} – are models in reality acting differently safe depending on user demographics?
Those two parts are, of course, interdependent, but the latter cannot be measured before we have trustworthy instruments, in this case demographically diverse safety evaluations, to do so.

The contributions of this thesis are the following four\footnote{The code for all empirical components of this thesis is available at \url{https://github.com/theaLilott/safe-for-whom}.}:
\begin{enumerate}
    \item A mapping of demographic diversity in evaluations and benchmarks concerned with user safety to establish the ground of methodological safety discrimination in safety evaluations (Chapter \ref{chap:evaleval}).
    \item A case study and discussion to enhance the linguistic and dialectal diversity (Chapter \ref{chap:translations}).
    \item Three case studies on demographic contextualisation at the judge-, prompt- and memory-level of safety evaluations (Chapter \ref{chap:contextualising}).
    \item A case study on cross-lingual generalisation of linear probes used for monitoring unwanted behaviour of LLMs (Chapter \ref{chap:probes}).
\end{enumerate}

Through this work, we aim to both guide future development of evaluations that show demographic diversity and enable measurements of behavioural safety discrimination. As more users worldwide use AI tools for a variety of tasks in their personal and professional lives, it is unquestionably important to ensure safeguards hold for all users and that AI leaves no one behind.
